\chapter{Apêndice ou Anexo?}
\label{ap:diferenca}

O Apêndice~\ref{ap:diferenca} explica a diferença entre apêndice e
anexo, conforme as convenções do \emph{template} oficial da
\thesisuniversity.

\paragraph{Apêndice.} Apêndices englobam materiais elaborados pelo
autor(a), tais como gráficos, quadros, tabelas, traduções,
organogramas e esquemas que prestem informação relevante para a
compreensão do trabalho. Só devem figurar nos apêndices as
informações previamente referenciadas no texto. As informações são
total ou parcialmente da responsabilidade do autor.

\paragraph{Anexo.} Anexos englobam documentos que, não sendo
elaborados pelo autor, serviram de base para a construção do estudo
ou facilitam a compreensão da tese/dissertação. Só devem figurar
nos anexos documentos e/ou materiais previamente referenciados no
corpo do trabalho. Todos os documentos devem estar em formato
digital.

% ============================================================================
\chapter{Cálculo de Spearman e Kendall em Python puro}
\label{ap:correlation-code}

Este apêndice apresenta a implementação completa do \emph{script}
\code{compute\_rank\_correlations.py} (Capítulo~\ref{ch:metodologia},
Secção~\ref{sec:metricas}), em Python puro e sem dependência de
\code{scipy}. O \emph{script} consome os ficheiros
\code{ranking\_results.json} produzidos pelo \code{ranker.py} e
gera o ficheiro consolidado \code{rank\_correlations.json},
utilizado para popular a Tabela~\ref{tab:correlacoes}.

\begin{lstlisting}[language=Python, caption={Implementação de Spearman $\rho$ e Kendall $\tau_b$ em Python puro (excerto).},label={lst:rank-code},basicstyle=\ttfamily\scriptsize]
def mean_ranks(scores):
    """Converte scores em ranks com mean-rank para ties."""
    indexed = sorted(enumerate(scores), key=lambda x: -x[1])
    ranks = [0.0] * len(scores)
    i = 0
    while i < len(indexed):
        j = i
        while j + 1 < len(indexed) and indexed[j+1][1] == indexed[i][1]:
            j += 1
        avg = (i + 1 + j + 1) / 2.0
        for k in range(i, j + 1):
            ranks[indexed[k][0]] = avg
        i = j + 1
    return ranks

def spearman_rho(a, b):
    return pearson(mean_ranks(a), mean_ranks(b))

def kendall_tau_b(a, b):
    """Kendall tau-b com correcao para empates (Goodman, 1954)."""
    n = len(a); C = D = Ta = Tb = 0
    for i in range(n):
        for j in range(i + 1, n):
            da, db = a[i] - a[j], b[i] - b[j]
            if   da == 0 and db == 0: Ta += 1; Tb += 1
            elif da == 0:             Ta += 1
            elif db == 0:             Tb += 1
            elif (da > 0) == (db > 0): C += 1
            else:                      D += 1
    P0 = n * (n - 1) / 2
    den = ((P0 - Ta) * (P0 - Tb)) ** 0.5
    return (C - D) / den if den else float('nan')
\end{lstlisting}

A versão completa, incluindo o carregamento dos ficheiros JSON e a
escrita do resultado consolidado, encontra-se em
\code{compute\_rank\_correlations.py} na raiz do repositório
público.

% ============================================================================
\chapter{Mapeamentos \code{func\_to\_req} dos três casos de estudo}
\label{ap:mapeamentos}

Os mapeamentos \code{func\_to\_req} utilizados em cada caso de
estudo encontram-se em:

\begin{itemize}
    \item CPython \emph{stdlib}: \code{testes/simples\_stdlib/mapeamento.py}
    \item Flask: \code{testes/medio\_flask/mapeamento.py}
    \item \emph{scikit-learn}: \code{testes/complexo\_sklearn/mapeamento.py}
\end{itemize}

Os ficheiros são distribuídos como parte do repositório público
mencionado em Capítulo~\ref{ch:introducao},
Secção~\ref{sec:contribuicoes}. Cada ficheiro contém um único
dicionário Python \code{FUNC\_TO\_REQ} cujas chaves são nomes
qualificados de funções/métodos e cujos valores são identificadores
de requisitos no formato \texttt{REQ\_*}.

% ============================================================================
\chapter{Representações Visuais do Call Graph}
\label{ap:callgraph-repr}

Este apêndice complementa a Figura~\ref{fig:call-graph-stdlib}
(Capítulo~\ref{ch:metodologia}) apresentando, para cada um dos três
casos de estudo, o \emph{Call Graph} extraído por análise estática
(Secção~\ref{sec:reqgraph}) em \textbf{três representações
complementares}. A motivação é metodológica: um único \emph{Call
Graph} real combina centenas de vértices com nomes longos de
funções/métodos, pelo que nenhuma representação isolada consegue, em
simultâneo, documentar a \emph{escala} do grafo e manter os
\emph{rótulos legíveis}. As três vistas resolvem este compromisso de
formas distintas:

\begin{description}
    \item[(i) Grafo completo.] Mostra \emph{todos} os vértices e
        arestas, dispostos por camadas hierárquicas (funções
        ``fonte'' no topo; folhas em baixo). Documenta fielmente a
        \textbf{escala e a densidade} reais do grafo. Os nomes estão
        rotulados em todos os vértices, mas, dada a dimensão, servem
        sobretudo de referência de conjunto --- não se destinam a
        leitura individual.
    \item[(ii) Maior componente conexa.] Isola a maior
        \emph{componente fracamente conexa} do grafo, eliminando os
        sub-sistemas periféricos e os vértices isolados. Revela a
        \textbf{estrutura principal} de invocações com rótulos já
        substancialmente mais legíveis.
    \item[(iii) Excerto representativo.] Um excerto conexo de
        $\sim\!10$ vértices, selecionado de forma determinística a
        partir do vértice de maior grau da maior componente
        (expansão em largura). Apresentado no estilo limpo dos grafos
        do Capítulo~\ref{ch:fundamentos} --- setas escuras e
        proeminentes que terminam na borda do nó ---, é a vista
        \textbf{totalmente legível}, adequada à inspeção dos nomes
        das funções/métodos em cada aresta.
\end{description}

As três figuras de cada caso foram geradas pelo \emph{script}
reprodutível \code{scripts/gen\_callgraph\_figures.py}, que reaplica
o extrator de \emph{Call Graph} (\code{reqgraph.call\_graph}) sobre os
ficheiros de cada caso de estudo. A cor de cada vértice identifica o
\emph{módulo} (ficheiro \code{.py}) de origem da função.

% ----------------------------------------------------------------------------
\section{Caso Simples --- CPython \emph{stdlib}}
\label{ap:cg-stdlib}

O \emph{Call Graph} do caso \emph{stdlib} (\code{argparse.py} +
\code{http/server.py}) tem $120$ vértices e $130$ arestas, repartidos
por vários sub-sistemas independentes --- coerente com o baixo
acoplamento esperado da biblioteca padrão. A
Figura~\ref{fig:cg-stdlib-full} evidencia esta dispersão em camadas; a
maior componente (Figura~\ref{fig:cg-stdlib-comp}) corresponde ao
\emph{core} de \emph{parsing} do \code{argparse}; o excerto
(Figura~\ref{fig:cg-stdlib-exc}) centra-se em
\code{ArgumentParser.\_parse\_known\_args}, o método que orquestra a
análise de argumentos.

\begin{figure}[htbp]
    \centering
    \includegraphics[width=\linewidth]{call_graph_simples_stdlib_full.png}
    \caption{CPython \emph{stdlib} --- \emph{Call Graph} completo
    ($120$ vértices, $130$ arestas).}
    \label{fig:cg-stdlib-full}
\end{figure}

\begin{figure}[htbp]
    \centering
    \includegraphics[width=0.92\linewidth]{call_graph_simples_stdlib_componente.png}
    \caption{CPython \emph{stdlib} --- maior componente conexa.}
    \label{fig:cg-stdlib-comp}
\end{figure}

\begin{figure}[htbp]
    \centering
    \includegraphics[width=0.85\linewidth]{call_graph_simples_stdlib_excerto.png}
    \caption{CPython \emph{stdlib} --- excerto representativo, no
    estilo do Capítulo~\ref{ch:fundamentos}.}
    \label{fig:cg-stdlib-exc}
\end{figure}

\clearpage
% ----------------------------------------------------------------------------
\section{Caso Intermediário --- Flask}
\label{ap:cg-flask}

O \emph{Call Graph} do Flask (\code{app.py}, \code{cli.py},
\code{blueprints.py}) tem $49$ vértices e $47$ arestas. A maior
componente (Figura~\ref{fig:cg-flask-comp}) converge para
\code{Flask.ensure\_sync} e para o ciclo de vida da requisição; o
excerto (Figura~\ref{fig:cg-flask-exc}) torna explícito o padrão em
estrela em torno de \code{Flask.ensure\_sync}, refletindo o seu papel
de ponto de convergência do \emph{middleware}.

\begin{figure}[htbp]
    \centering
    \includegraphics[width=\linewidth]{call_graph_medio_flask_full.png}
    \caption{Flask --- \emph{Call Graph} completo ($49$ vértices,
    $47$ arestas).}
    \label{fig:cg-flask-full}
\end{figure}

\begin{figure}[htbp]
    \centering
    \includegraphics[width=0.92\linewidth]{call_graph_medio_flask_componente.png}
    \caption{Flask --- maior componente conexa.}
    \label{fig:cg-flask-comp}
\end{figure}

\begin{figure}[htbp]
    \centering
    \includegraphics[width=0.85\linewidth]{call_graph_medio_flask_excerto.png}
    \caption{Flask --- excerto representativo, no estilo do
    Capítulo~\ref{ch:fundamentos}.}
    \label{fig:cg-flask-exc}
\end{figure}

\clearpage
% ----------------------------------------------------------------------------
\section{Caso Complexo --- \emph{scikit-learn}}
\label{ap:cg-sklearn}

O \emph{Call Graph} do \emph{scikit-learn} (\code{pipeline.py},
\code{linear\_model/\_base.py}, \code{linear\_model/\_logistic.py},
\code{tree/\_classes.py}) tem $74$ vértices e $80$ arestas. A maior
componente (Figura~\ref{fig:cg-sklearn-comp}) e o excerto
(Figura~\ref{fig:cg-sklearn-exc}) confirmam o papel central de
\code{Pipeline.\_iter} e dos métodos \code{fit}/\code{predict} como
orquestradores --- a contraparte, ao nível das funções, da
centralidade de \req{REQ\_PIPELINE} observada no Grafo de Requisitos
(Secção~\ref{sec:res-sklearn}).

\begin{figure}[htbp]
    \centering
    \includegraphics[width=\linewidth]{call_graph_complexo_sklearn_full.png}
    \caption{\emph{scikit-learn} --- \emph{Call Graph} completo
    ($74$ vértices, $80$ arestas).}
    \label{fig:cg-sklearn-full}
\end{figure}

\begin{figure}[htbp]
    \centering
    \includegraphics[width=0.92\linewidth]{call_graph_complexo_sklearn_componente.png}
    \caption{\emph{scikit-learn} --- maior componente conexa.}
    \label{fig:cg-sklearn-comp}
\end{figure}

\begin{figure}[htbp]
    \centering
    \includegraphics[width=0.85\linewidth]{call_graph_complexo_sklearn_excerto.png}
    \caption{\emph{scikit-learn} --- excerto representativo, no estilo
    do Capítulo~\ref{ch:fundamentos}.}
    \label{fig:cg-sklearn-exc}
\end{figure}
